\subsection{Introduction: The Data Frontier \quad / \quad \textthai{บทนำ: พรมแดนแห่งข้อมูล}}
\begin{paracol}{2}
We conclude our journey at the extreme frontier of physics: High-Energy Physics (HEP)\index{High-Energy Physics@ฟิสิกส์พลังงานสูง (High-Energy Physics)}. In facilities like the Large Hadron Collider (LHC), subatomic particles collide at nearly the speed of light, generating petabytes of raw data per second. This "Data Tsunami" is impossible to store in its entirety. Instead, physicists must use automated \textbf{Trigger Systems}\index{Triggers@ทริกเกอร์ (Triggers)}—intelligent filters that decide in real-time which 0.001\% of events are worth keeping.

In this chapter, we explore how AI has moved from a simple filter to an active scientific collaborator. We will focus on \textbf{Graph Neural Networks (GNNs)}\index{GNNs@โครงข่ายประสาทเทียมแบบกราฟ (GNNs)}, a technology that treats detector hits as nodes in a graph, allowing for the precise reconstruction of particle tracks. Finally, we discuss the \textbf{Ethics of Discovery}\index{Ethics@จริยธรรม (Ethics)}: how we can trust AI when it discoveries new physics that humanity cannot yet explain.

\switchcolumn

\begin{thai}
เราสรุปการเดินทางของเราที่พรมแดนปลายสุดของฟิสิกส์: ฟิสิกส์พลังงานสูง (HEP) ในสิ่งอำนวยความสะดวกอย่าง Large Hadron Collider (LHC) อนุภาคระดับต่ำกว่าอะตอมจะพุ่งชนกันด้วยความเร็วเกือบเท่าแสง สร้างข้อมูลดิบระดับเพตาไบต์ต่อวินาที "สึนามิข้อมูล" นี้เป็นไปไม่ได้ที่จะจัดเก็บทั้งหมด นักฟิสิกส์จึงต้องใช้ \textbf{ระบบทริกเกอร์ (Trigger Systems)} อัตโนมัติ ซึ่งเป็นตัวกรองอัจฉริยะที่ตัดสินใจแบบเรียลไลม์ว่า 0.001\% ของเหตุการณ์ใดที่คุ้มค่าแก่การเก็บรักษา

ในบทนี้ เราจะสำรวจว่า AI เปลี่ยนจากตัวกรองธรรมดาไปสู่การเป็นผู้ร่วมงานทางวิทยาศาสตร์ที่กระตือรือร้นได้อย่างไร เราจะมุ่งเน้นไปที่ \textbf{Graph Neural Networks (GNNs)} ซึ่งเป็นเทคโนโลยีที่มองว่าการตรวจจับอนุภาคเป็นโหนดในกราฟ ช่วยให้สามารถสร้างเส้นทางอนุภาคใหม่ได้อย่างแม่นยำ สุดท้าย เราจะหารือเกี่ยวกับ \textbf{จริยธรรมแห่งการค้นพบ (Ethics of Discovery)} ว่าเราจะเชื่อถือ AI ได้อย่างไรเมื่อมันค้นพบฟิสิกส์ใหม่ๆ ที่มนุษยชาติยังไม่สามารถอธิบายได้
\end{thai}
\end{paracol}
